% ==============================================================================
%  MDA · Tema 4: SCRUM — apuntes y ejercicios para el parcial del 27/11/2026
%  Fuente: diapositivas «Unidad 1 - Desarrollo Ágil de Software. Tema 4 – SCRUM»
%  (104 diapositivas) y «Tema 4 – Elevator Pitch» (11 diapositivas), LSI, UGR.
%  Compilar: pdflatex tema4_scrum.tex (dos pasadas)
% ==============================================================================
\documentclass[11pt,a4paper,oneside,openany]{report}
% ==============================================================================
% CONFIGURACIÓN Y ESTILO VISUAL — METODOLOGÍAS DE DESARROLLO ÁGIL (MDA)
% Departamento de Lenguajes y Sistemas Informáticos — Universidad de Granada
% Plataforma repasaYA
% ==============================================================================

\usepackage[utf8]{inputenc}
\usepackage[T1]{fontenc}
\usepackage[spanish,es-tabla]{babel}
\usepackage{lmodern}
\usepackage{microtype}
\usepackage[dvipsnames,svgnames,table,xcdraw]{xcolor}
\usepackage{amsmath,amssymb,amsfonts,amsthm}
\usepackage{graphicx}
\usepackage{float}
\usepackage{booktabs}
\usepackage{tabularx}
\usepackage{longtable}
\usepackage{array}
\usepackage{multirow}
\usepackage{multicol}
\usepackage{eurosym}
\usepackage{enumitem}
\usepackage{fancyhdr}
\usepackage{titlesec}
\usepackage{listings}
\usepackage{tcolorbox}
\tcbuselibrary{skins,breakable}
\usepackage{tikz}
\usetikzlibrary{babel,arrows.meta,positioning,calc,shapes.geometric,fit,backgrounds}

% Configuración de márgenes
\usepackage[a4paper,top=2.5cm,bottom=2.5cm,left=2.6cm,right=2.6cm,headheight=24pt]{geometry}

% Paleta de colores profesionales para Ingeniería del Software / Ágil
\definecolor{mdaDarkBlue}{RGB}{18, 32, 58}
\definecolor{mdaPrimary}{RGB}{30, 74, 138}
\definecolor{mdaAccent}{RGB}{99, 91, 255}
\definecolor{mdaLightAccent}{RGB}{242, 243, 255}
\definecolor{mdaGreen}{RGB}{22, 135, 70}
\definecolor{mdaLightGreen}{RGB}{240, 253, 244}
\definecolor{mdaOrange}{RGB}{217, 119, 6}
\definecolor{mdaLightOrange}{RGB}{255, 251, 235}
\definecolor{mdaRed}{RGB}{185, 28, 28}
\definecolor{mdaLightRed}{RGB}{254, 242, 242}
\definecolor{mdaGray}{RGB}{246, 248, 251}
\definecolor{mdaBorderGray}{RGB}{215, 222, 232}
\definecolor{mdaTeal}{RGB}{15, 118, 110}
\definecolor{mdaYellowCard}{RGB}{255, 252, 220}
\definecolor{mdaYellowBorder}{RGB}{210, 185, 80}
\definecolor{codeBg}{RGB}{248, 249, 252}

% Hipervínculos
\usepackage{hyperref}
\hypersetup{
    colorlinks=true,
    linkcolor=mdaDarkBlue,
    citecolor=mdaTeal,
    urlcolor=mdaAccent,
    pdfauthor={repasaYA - Universidad de Granada (UGR)},
    pdftitle={Metodologías de Desarrollo Ágil - Seminario 2: Historias de Usuario},
    pdfsubject={Desarrollo Ágil de Software, Scrum, XP, Historias de Usuario, INVEST, Planning Poker, Spikes},
    pdfkeywords={MDA, Historias de Usuario, Scrum, XP, INVEST, CCC, BDD, Planning Poker, Velocidad, Spikes, UGR}
}

% Configuración de Encabezados y Pies de página
\pagestyle{fancy}
\fancyhf{}
\fancyhead[LE,RO]{\footnotesize\textsf{\thepage}}
\fancyhead[LO]{\footnotesize\textsf{\rightmark}}
\fancyhead[RE]{\footnotesize\textsf{\leftmark}}
\fancyfoot[C]{\footnotesize\textsf{\color{gray}repasaYA $\cdot$ Metodologías de Desarrollo Ágil (UGR) $\cdot$ Seminario 2: Historias de Usuario}}
\renewcommand{\headrulewidth}{0.4pt}
\renewcommand{\footrulewidth}{0.2pt}

% Configuración de títulos de capítulos y secciones
\titleformat{\chapter}[display]
  {\normalfont\bfseries\color{mdaDarkBlue}\sffamily}
  {\Large\textsf{Bloque \thechapter}}{1ex}
  {\titlerule[1.5pt]\vspace{1ex}\huge}
  [\vspace{1ex}\titlerule]

\titleformat{\section}
  {\normalfont\Large\bfseries\color{mdaPrimary}\sffamily}
  {\thesection}{0.8em}{}

\titleformat{\subsection}
  {\normalfont\large\bfseries\color{mdaDarkBlue}\sffamily}
  {\thesubsection}{0.8em}{}

\titleformat{\subsubsection}
  {\normalfont\normalsize\bfseries\color{mdaAccent}\sffamily}
  {\thesubsubsection}{0.8em}{}

% Entorno: Caja de Fuentes al inicio de cada tema
\newcommand{\cuadrofuentes}[1]{%
  \begin{tcolorbox}[
    colback=mdaLightAccent,
    colframe=mdaAccent,
    boxrule=1pt,
    arc=4pt,
    left=12pt,
    right=12pt,
    top=8pt,
    bottom=8pt,
    breakable
  ]
    \small
    \noindent\textbf{\color{mdaAccent}\large\textsf{Síntesis del Bloque y Objetivos de Evaluación:}}\\[4pt]
    #1
  \end{tcolorbox}
  \vspace{8pt}
}

% Entorno: Definición
\newtcolorbox{definicion}[1][]{
  enhanced,
  colback=mdaGray,
  colframe=mdaPrimary,
  boxrule=1pt,
  leftrule=4pt,
  arc=3pt,
  title=\textbf{\textsf{Definición Canónica\ifstrempty{#1}{}{: #1}}},
  coltitle=white,
  colbacktitle=mdaPrimary,
  attach boxed title to top left={yshift=-2mm, xshift=4mm},
  boxed title style={arc=2pt, boxrule=0pt},
  top=10pt,
  bottom=8pt,
  left=10pt,
  right=10pt,
  breakable
}

% Entorno: Ejemplo Explicado
\newtcolorbox{ejemplo}[1][]{
  enhanced,
  colback=mdaLightGreen,
  colframe=mdaGreen,
  boxrule=0.8pt,
  leftrule=3.5pt,
  arc=3pt,
  title=\textbf{\textsf{Ejemplo de Clase / Dominio\ifstrempty{#1}{}{: #1}}},
  coltitle=white,
  colbacktitle=mdaGreen,
  attach boxed title to top left={yshift=-2mm, xshift=4mm},
  boxed title style={arc=2pt, boxrule=0pt},
  top=10pt,
  bottom=8pt,
  left=10pt,
  right=10pt,
  breakable
}

% Entorno: Aviso / Ojo Examen
\newtcolorbox{aviso}[1][]{
  enhanced,
  colback=mdaLightOrange,
  colframe=mdaOrange,
  boxrule=0.8pt,
  leftrule=3.5pt,
  arc=3pt,
  title=\textbf{\textsf{Importante / Pregunta de Examen\ifstrempty{#1}{}{: #1}}},
  coltitle=white,
  colbacktitle=mdaOrange,
  attach boxed title to top left={yshift=-2mm, xshift=4mm},
  boxed title style={arc=2pt, boxrule=0pt},
  top=10pt,
  bottom=8pt,
  left=10pt,
  right=10pt,
  breakable
}

% Entorno: Tarjeta de Historia de Usuario
\newtcolorbox{tarjetahu}[1][]{
  enhanced,
  colback=mdaYellowCard,
  colframe=mdaYellowBorder,
  boxrule=1pt,
  leftrule=4pt,
  arc=3pt,
  title=\textbf{\textsf{\color{mdaDarkBlue}Tarjeta de Historia de Usuario\ifstrempty{#1}{}{: #1}}},
  coltitle=mdaDarkBlue,
  colbacktitle=mdaYellowBorder!40!white,
  attach boxed title to top left={yshift=-2mm, xshift=4mm},
  boxed title style={arc=2pt, boxrule=0.8pt, colframe=mdaYellowBorder},
  top=10pt,
  bottom=8pt,
  left=10pt,
  right=10pt,
  breakable
}

% Entorno: Regla Formal / Patrón
\newtcolorbox{teorema}[1][]{
  enhanced,
  colback=mdaLightAccent,
  colframe=mdaAccent,
  boxrule=0.8pt,
  leftrule=3.5pt,
  arc=3pt,
  title=\textbf{\textsf{Principio / Patrón Metodológico\ifstrempty{#1}{}{: #1}}},
  coltitle=white,
  colbacktitle=mdaAccent,
  attach boxed title to top left={yshift=-2mm, xshift=4mm},
  boxed title style={arc=2pt, boxrule=0pt},
  top=10pt,
  bottom=8pt,
  left=10pt,
  right=10pt,
  breakable
}

% Entorno: Ejercicio Resuelto
\newtcolorbox{ejercicio}[1][]{
  enhanced,
  colback=white,
  colframe=mdaDarkBlue,
  boxrule=1pt,
  leftrule=4pt,
  arc=3pt,
  title=\textbf{\textsf{Ejercicio de Examen Resuelto\ifstrempty{#1}{}{: #1}}},
  coltitle=white,
  colbacktitle=mdaDarkBlue,
  attach boxed title to top left={yshift=-2mm, xshift=4mm},
  boxed title style={arc=2pt, boxrule=0pt},
  top=10pt,
  bottom=8pt,
  left=10pt,
  right=10pt,
  breakable
}

% Configuración de listings para Gherkin / BDD
\lstset{
  backgroundcolor=\color{codeBg},
  basicstyle=\ttfamily\small\color{black},
  keywordstyle=\bfseries\color{mdaPrimary},
  commentstyle=\itshape\color{gray},
  stringstyle=\color{mdaGreen},
  numbers=none,
  breaklines=true,
  breakatwhitespace=false,
  showstringspaces=false,
  frame=single,
  rulecolor=\color{mdaBorderGray},
  tabsize=2,
  captionpos=b,
  inputencoding=utf8,
  extendedchars=true,
  literate=
    {á}{{\'a}}1 {é}{{\'e}}1 {í}{{\'i}}1 {ó}{{\'o}}1 {ú}{{\'u}}1
    {Á}{{\'A}}1 {É}{{\'E}}1 {Í}{{\'I}}1 {Ó}{{\'O}}1 {Ú}{{\'U}}1
    {ñ}{{\~n}}1 {Ñ}{{\~N}}1 {¿}{{?`}}1 {¡}{{!`}}1
    {«}{{\guillemotleft}}1 {»}{{\guillemotright}}1
    {€}{{\euro}}1
}


\renewcommand{\thesection}{\arabic{section}}
\setcounter{secnumdepth}{2}
\newcommand{\diap}[1]{\textcolor{mdaPrimary}{\small[diap.~#1]}}
\newcommand{\extra}{\textcolor{mdaOrange}{\small[no está en las diapositivas]}}

\hypersetup{pdftitle={MDA - Tema 4: Scrum. Apuntes y ejercicios para el parcial}}
\fancyfoot[C]{\footnotesize\textsf{\color{gray}repasaYA $\cdot$ Metodologías de Desarrollo Ágil (UGR) $\cdot$ Tema 4: Scrum}}

\begin{document}

% ------------------------------------------------------------------------------
\begin{titlepage}
  \centering
  \vspace*{1.2cm}
  {\Huge\bfseries\color{mdaDarkBlue}\textsf{Metodologías de Desarrollo Ágil}}\\[0.5cm]
  {\LARGE\bfseries\color{mdaPrimary}\textsf{Tema 4: Scrum}}\\[0.3cm]
  {\Large\textsf{Apuntes, esquemas y ejercicios para el segundo parcial}}\\[1.2cm]
  \begin{tcolorbox}[colback=mdaLightAccent,colframe=mdaAccent,width=0.88\textwidth,arc=4pt,halign=center]
    \large\textbf{Viernes 27 de noviembre de 2026}\\[4pt]
    \normalsize Prueba \textbf{en grupo} (los grupos de prácticas) $\cdot$ \textbf{sin apuntes} $\cdot$ \textbf{1 hora como máximo}\\[4pt]
    Preguntas cortas y ejercicios sobre el contenido del tema
  \end{tcolorbox}
  \vspace{0.8cm}
  \begin{minipage}{0.88\textwidth}\small
  \textbf{De dónde sale.} De las 104 diapositivas del Tema 4 y de las 11 del \emph{Elevator Pitch} (Sprint 0 y
  visión del producto). Cada idea lleva la diapositiva de la que sale, por ejemplo \diap{53}. Lo poco que se añade de
  fuera (la duración del Daily, la plantilla de Moore para el \emph{elevator pitch}) se marca con
  \extra.\\[6pt]
  \textbf{Cómo usarlo.} Secciones 1 a 6: la teoría en el orden de las diapositivas. Sección 7: cómo repartiros el
  trabajo en el examen de grupo. Sección 8: ejercicios resueltos de los tipos que pueden caer. Sección 9: preguntas
  cortas con su respuesta modelo. Al final, una tabla resumen para el último repaso.
  \end{minipage}
  \vfill
  {\large\textsf{repasaYA $\cdot$ Universidad de Granada}}
\end{titlepage}

\tableofcontents
\clearpage

% ==============================================================================
\section{Qué es Scrum}

\subsection{Una nueva forma de desarrollar \diap{2}}
Takeuchi y Nonaka (\emph{The New New Product Development Game}, Harvard Business Review, 1986) comparan dos formas
de desarrollar productos:
\begin{itemize}[leftmargin=*]
  \item \textbf{Carrera de relevos}: cada fase pasa el testigo a la siguiente (como en cascada). Choca con la
        velocidad y la flexibilidad que exige la competencia.
  \item \textbf{Rugby}: el equipo avanza \emph{como una unidad}, pasándose el balón adelante y atrás. De ahí viene el
        nombre: un \emph{scrum} es la melé del rugby.
\end{itemize}
Por eso, en lugar de hacer todo de una cosa a la vez, \textbf{los equipos Scrum hacen un poco de todo, todo el
tiempo}: en cada Sprint se analiza, diseña, codifica y prueba \diap{48}.

\subsection{Historia \diap{3--5}}
\begin{center}\small
\begin{tabular}{lp{11.6cm}}
\toprule
\textbf{Año} & \textbf{Hito} \\ \midrule
1986 & Takeuchi y Nonaka usan por primera vez el término, para un proceso de desarrollo acelerado y competitivo. \\
1993 & Primer equipo que trabaja con Scrum: \textbf{Jeff Sutherland} en \textbf{Easel Corp.} \\
1996 & \textbf{Ken Schwaber} formaliza Scrum en el artículo \emph{Scrum Development Process} (conferencia \textbf{OOPSLA}). \\
2001 & Schwaber y \textbf{Mike Beedle} publican el primer libro: \emph{Agile Software Development with Scrum}. \\
2002 & Schwaber y \textbf{Mike Cohn} fundan la \textbf{Scrum Alliance}, para compartir y aumentar el conocimiento de Scrum. \\
-- & Schwaber y Sutherland mantienen \emph{La Guía de Scrum} (scrum.org). \\
\bottomrule
\end{tabular}
\end{center}
Se usa en empresas de todo tipo (Google, Microsoft, Amazon, Telefónica, BBC…) y en todo tipo de software:
videojuegos, sistemas críticos, contratos de precio cerrado, aplicaciones certificadas ISO 9001, sistemas
embebidos, sistemas 7x24 con 99,999~\% de disponibilidad… \diap{6--7}. Scrum \textbf{no} es solo para proyectos
pequeños o web.

\subsection{Definición \diap{8}}
\begin{definicion}[Scrum]
Scrum es un \textbf{marco de trabajo iterativo e incremental} para el desarrollo de proyectos, productos y
aplicaciones. \textbf{No se basa en el seguimiento de un plan}, sino en la \textbf{adaptación continua} a las
circunstancias de la evolución del proyecto.
\end{definicion}
\begin{itemize}[leftmargin=*]
  \item \textbf{Iterativo}: el trabajo se repite en ciclos (Sprints).
  \item \textbf{Incremental}: cada ciclo añade un trozo de producto terminado.
  \item Es un \textbf{marco} (\emph{framework}), no una metodología cerrada: dice qué roles, eventos y artefactos hay,
        no cómo programar.
\end{itemize}

\subsection{Pilares empíricos \diap{9}}
Scrum es \textbf{empírico}: se decide con lo que se observa, no con lo que se predijo.
\begin{center}
\begin{tikzpicture}[font=\sffamily\small]
  \foreach \x/\t/\d in {0/Transparencia/{Todos ven el\\estado real},4.6/Inspección/{Se revisa a menudo\\el producto y el proceso},9.2/Adaptación/{Se corrige en cuanto\\algo se desvía}}{
    \node[draw=mdaPrimary,fill=mdaLightAccent,rounded corners,minimum width=3.9cm,minimum height=1.6cm,align=center]
      at (\x,0) {\textbf{\t}\\[2pt]\footnotesize\d};}
  \draw[-{Latex},thick,mdaPrimary] (1.95,0) -- (2.65,0);
  \draw[-{Latex},thick,mdaPrimary] (6.55,0) -- (7.25,0);
\end{tikzpicture}
\end{center}
Ejemplos: el Burndown Chart y el panel de tareas dan \emph{transparencia}; el Daily y la Review son
\emph{inspección}; replanificar el Sprint o cambiar la forma de trabajar tras la Retrospectiva es
\emph{adaptación}.

\subsection{Valores \diap{10}}
Las diapositivas listan catorce: \textbf{mejora continua, calidad, time-boxing, responsabilidad,
multidisciplinar, flexibilidad, ritmo, compromiso, simplicidad, respeto, coraje, foco, predictibilidad y personas}.
\begin{aviso}[los cinco valores de la Guía de Scrum]
La Guía de Scrum oficial reduce los valores a cinco: \textbf{compromiso, coraje, foco, apertura y respeto}
\extra. Si preguntan «valores de Scrum», da los de la diapositiva y, si quieres, menciona estos cinco.
\end{aviso}

\subsection{La esencia de Scrum \diap{11--12} y factores clave \diap{103}}
\begin{multicols}{2}\small
\textbf{Esencia}
\begin{itemize}[leftmargin=*,itemsep=1pt]
  \item El equipo recibe \textbf{objetivos claros}.
  \item Se \textbf{organiza} en función del trabajo.
  \item Entrega con \textbf{regularidad} las funciones \textbf{más valiosas}.
  \item Recibe \textbf{retroalimentación} de gente de fuera del equipo.
  \item \textbf{Reflexiona} sobre su forma de trabajar para mejorar.
  \item Toda la organización tiene \textbf{visibilidad} del progreso.
  \item Equipo y gerencia se comunican de forma \textbf{honesta y transparente}.
\end{itemize}
\columnbreak
\textbf{Factores clave}
\begin{itemize}[leftmargin=*,itemsep=1pt]
  \item \textbf{Delegación} de atribuciones al equipo: auto-organización y toma de decisiones.
  \item \textbf{Respeto} entre las personas: confianza en sus conocimientos y capacidades.
  \item \textbf{Responsabilidad} y autodisciplina.
  \item Trabajo centrado en el \textbf{compromiso} de desarrollo.
  \item \textbf{Información, transparencia y visibilidad}.
\end{itemize}
\end{multicols}

% ==============================================================================
\section{Elementos de Scrum \diap{13}}
\begin{center}
\begin{tikzpicture}[font=\sffamily\small,
  cab/.style={draw=mdaPrimary,fill=mdaPrimary,text=white,rounded corners,minimum width=4.2cm,minimum height=0.7cm,font=\sffamily\bfseries},
  it/.style={draw=mdaBorderGray,fill=mdaLightAccent,minimum width=4.2cm,minimum height=0.85cm,align=center,font=\sffamily\footnotesize}]
  \node[cab] at (0,0) {Artefactos (3)};
  \node[it] at (0,-0.9) {Product Backlog\\(Pila de producto)};
  \node[it] at (0,-1.8) {Sprint Backlog\\(Pila de Sprint)};
  \node[it] at (0,-2.7) {Burndown Chart\\(Gráfica de trabajo pendiente)};
  \node[cab] at (4.7,0) {Reuniones (4)};
  \node[it] at (4.7,-0.9) {Sprint Planning\\(Planificación del Sprint)};
  \node[it] at (4.7,-1.8) {Daily Meeting\\(Reunión diaria)};
  \node[it] at (4.7,-2.7) {Sprint Review\\(Revisión del Sprint)};
  \node[it] at (4.7,-3.6) {Sprint Retrospective\\(Retrospectiva)};
  \node[cab] at (9.4,0) {Roles (3)};
  \node[it] at (9.4,-0.9) {Product Owner\\(Propietario)};
  \node[it] at (9.4,-1.8) {Scrum Master\\(Director de proyecto)};
  \node[it] at (9.4,-2.7) {Team\\(Equipo de desarrollo)};
  \node[text=mdaRed,font=\sffamily\bfseries] at (4.7,-4.5) {+ Reglas de trabajo};
\end{tikzpicture}
\end{center}
\textbf{Regla mnemotécnica: 3 -- 4 -- 3} (tres artefactos, cuatro reuniones, tres roles), más las reglas de
trabajo.

% ==============================================================================
\section{Roles}

\subsection{Rol $\neq$ puesto \diap{15}}
Los roles representan una \textbf{responsabilidad en el proceso}, \textbf{no la posición} dentro de la organización.
Un jefe de departamento puede ser un miembro más del equipo, y un programador puede hacer de Scrum Master.

\subsection{Gallinas y cerdos: 3 + 3 roles \diap{16}}
\begin{center}
\begin{tikzpicture}[font=\sffamily\small]
  \node[draw=mdaOrange,fill=mdaLightOrange,rounded corners,minimum width=6cm,minimum height=1.1cm,align=center] (g) at (0,0)
    {\textbf{Roles gallina} (stakeholders)\\\footnotesize User $\cdot$ Customer $\cdot$ Management};
  \node[draw=mdaPrimary,fill=mdaLightAccent,rounded corners,minimum width=6cm,minimum height=1.1cm,align=center] (c) at (0,-1.8)
    {\textbf{Roles cerdo} (equipo Scrum)\\\footnotesize Product Owner $\cdot$ Scrum Master $\cdot$ Team};
  \draw[dashed,mdaOrange] (-3.6,-0.9) -- (3.6,-0.9);
  \node[align=left,font=\sffamily\footnotesize,text width=6.2cm] at (7.2,-0.9) {
    En unos huevos con beicon, la gallina \emph{participa} (pone los huevos) y el cerdo se \emph{compromete} (pone
    el beicon). Los cerdos se juegan el resultado; las gallinas están interesadas, pero no comprometidas.};
\end{tikzpicture}
\end{center}
\textbf{Stakeholders} \diap{17}: todas las personas y organizaciones con algún interés en el trabajo. Son
\textbf{fuente de requisitos} (definen el producto) y \textbf{validan el trabajo} (examinan los resultados intermedios y
finales y dan su \emph{feedback}).

\subsection{Product Owner (propietario) \diap{18--21}}
\begin{definicion}[Product Owner]
\textbf{Representante del cliente e intermediario entre éste y el equipo.} Habla el lenguaje del negocio y conoce
los métodos y conceptos del equipo. Es \textbf{una sola persona}, no un comité.
\end{definicion}
\textbf{Implementación} \diap{19}: \emph{directa} (el propio cliente hace de PO) o \emph{proxy} (alguien de la
organización, como un \emph{product manager}, representa a varios clientes ante el equipo).

\begin{multicols}{2}\small
\textbf{Necesidades del PO} \diap{20}
\begin{itemize}[leftmargin=*,itemsep=1pt]
  \item Conocer a fondo el negocio del cliente, sus necesidades y el objetivo del sistema.
  \item Tener \textbf{atribuciones} para tomar decisiones.
  \item Conocer Scrum.
  \item Recibir y analizar retroinformación continua del negocio y del proyecto.
\end{itemize}
\columnbreak
\textbf{Responsabilidades del PO} \diap{21}
\begin{itemize}[leftmargin=*,itemsep=1pt]
  \item Obtener una \textbf{visión compartida}.
  \item Recoger los requisitos.
  \item Crear, administrar y \textbf{priorizar el Product Backlog}.
  \item \textbf{Aceptar o rechazar} el software al final de cada iteración.
  \item Resolver las dudas del equipo.
  \item Administrar el \textbf{plan de entregas}.
  \item \textbf{Maximizar el ROI}.
\end{itemize}
\end{multicols}

\subsection{Scrum Master \diap{22--28}}
\begin{definicion}[Scrum Master]
Garantiza el funcionamiento de los procesos y metodologías que se emplean, gestiona el proceso y mejora la
productividad del equipo. \textbf{Es un facilitador, coach y mentor}, no un jefe: \textbf{no asigna tareas} ni
manda sobre el equipo.
\end{definicion}
\begin{aviso}[trampa de nombre]
Las diapositivas traducen Scrum Master como «director de proyecto», pero \textbf{no es un jefe de proyecto
tradicional}. No reparte el trabajo (el equipo se auto-organiza) ni decide qué se construye (eso es del PO).
\end{aviso}
\begin{multicols}{2}\small
\textbf{Responsabilidades} \diap{23}
\begin{itemize}[leftmargin=*,itemsep=1pt]
  \item Promover los valores y prácticas de Scrum.
  \item \textbf{Eliminar impedimentos}.
  \item Asegurar que el equipo es funcional y productivo.
  \item Permitir la cooperación entre todos los roles.
  \item \textbf{Escudo} del equipo ante interferencias externas.
  \item Fomentar el uso y el respeto del proceso.
\end{itemize}
\textbf{Atributos} \diap{24}: responsable, humilde, colaborador, comprometido, influyente y entendido.
\columnbreak
\textbf{Servicio al PO} \diap{25}: técnicas para gestionar la Lista de Producto; entender la planificación en un
entorno empírico; enseñarle a ordenar la lista para maximizar el valor.\\[4pt]
\textbf{Servicio al equipo} \diap{26}: guiarlo a ser auto-organizado y multifuncional; ayudarle a crear productos de
alto valor; eliminar impedimentos; facilitar los eventos; enseñarle Scrum donde la organización aún no lo entiende.\\[4pt]
\textbf{Servicio a la organización} \diap{27}: liderar la adopción de Scrum; planificar sus implantaciones; ayudar a
entender el desarrollo empírico; motivar cambios que suban la productividad; coordinarse con otros Scrum Masters.
\end{multicols}

\subsection{Equipo de desarrollo (Team) \diap{29--33}}
\begin{definicion}[Equipo de desarrollo]
Los profesionales que entregan, al final de cada Sprint, un \textbf{incremento de producto «terminado»} que
potencialmente se puede poner en producción.
\end{definicion}
\begin{itemize}[leftmargin=*]
  \item \textbf{Auto-organizado}: tiene la responsabilidad final del éxito y \textbf{no está jerarquizado}.
  \item \textbf{Multidisciplinar}: entre todos tienen las habilidades necesarias.
  \item \textbf{Multifuncional}: programar, diseñar, analizar, probar…
  \item Elevado grado de \textbf{compromiso}.
  \item Scrum \textbf{no reconoce sub-equipos} (nada de «equipo de pruebas» dentro del equipo).
  \item Trabajo a \textbf{tiempo completo} y de forma \textbf{fija}.
  \item \textbf{Tamaño}: lo bastante pequeño para seguir siendo ágil y lo bastante grande para completar trabajo
        significativo: \textbf{típicamente de 5 a 9 personas} \diap{31}.
  \item Cualidades: trabajo en equipo, comunicación, generosidad y capacidad de aprendizaje \diap{29}.
\end{itemize}
\textbf{Responsabilidades} \diap{32} (del equipo \emph{como un todo}): \textbf{estimar} el tamaño de los ítems del
backlog, comprometerse a entregar incrementos de \textbf{calidad de producción} y hacer el seguimiento de \textbf{su
propio plan}. Todos conocen la visión, colaboran con el PO en la pila, comparten el objetivo del Sprint, participan en
las decisiones y se respetan \diap{33}.

\subsection{¿Quién hace qué?}
\begin{center}\small
\begin{tabular}{p{6.3cm}ccc}
\toprule
\textbf{Tarea} & \textbf{PO} & \textbf{SM} & \textbf{Equipo} \\ \midrule
Definir la visión del producto & \textbf{X} & & \\
Crear, mantener y \textbf{priorizar} el Product Backlog & \textbf{X} & & colabora \\
\textbf{Estimar} historias (puntos de historia) & & & \textbf{X} \\
Descomponer historias en tareas (Sprint Backlog) & & & \textbf{X} \\
Elegir \emph{cómo} hacer el trabajo / repartirse las tareas & & & \textbf{X} \\
Aceptar o rechazar el incremento & \textbf{X} & & \\
Eliminar impedimentos; mantener la lista de impedimentos & & \textbf{X} & \\
Facilitar las reuniones y proteger al equipo & & \textbf{X} & \\
Maximizar el ROI y gestionar el plan de entregas & \textbf{X} & & \\
Actualizar el Burndown y el estado de las tareas & & & \textbf{X} \\
\bottomrule
\end{tabular}
\end{center}

% ==============================================================================
\section{Artefactos}
Los backlogs son los repositorios donde PO, SM y equipo escriben los requisitos (con \textbf{historias de usuario})
y las tareas de desarrollo \diap{35}.

\subsection{Product Backlog (pila de producto) \diap{36--39}}
\begin{definicion}[Product Backlog]
«Lista \textbf{ordenada y priorizada} de todo lo que podría ser necesario en el producto». Es \textbf{flexible},
\textbf{adaptable} y \textbf{dinámica}, está escrita en \textbf{lenguaje de negocio} y su responsable es el
\textbf{Product Owner}.
\end{definicion}
\begin{itemize}[leftmargin=*]
  \item Describe los requisitos que se conocen \textbf{en un momento dado}: nunca está completa.
  \item Los requisitos se discuten, aclaran y estiman entre todo el equipo Scrum, de forma colectiva.
  \item Contiene requisitos \textbf{funcionales y no funcionales}.
  \item Es la \textbf{materialización de la visión} y el nexo de comunicación entre clientes y equipo \diap{102}.
\end{itemize}
\begin{teorema}[Lista DEEP \diap{38}]
\begin{description}[leftmargin=1.2cm,style=nextline]
  \item[\textbf{D}etallado] Las HU están definidas de forma que el equipo pueda entenderlas y estimarlas (más detalle
        arriba, en las prioritarias; menos abajo).
  \item[\textbf{E}mergente] En cualquier momento pueden aparecer nuevas historias.
  \item[\textbf{E}stimado] Los elementos están estimados.
  \item[\textbf{P}riorizado] Para saber qué hay que detallar y dónde colocar los ítems nuevos.
\end{description}
\end{teorema}
\textbf{Técnicas de priorización} \diap{39} (Mike Cohn, 2008):
\begin{itemize}[leftmargin=*,itemsep=1pt]
  \item \textbf{Priority poker}: como el Planning Poker, pero votando prioridad.
  \item \textbf{MoSCoW}: \emph{Must} (imprescindible), \emph{Should} (importante), \emph{Could} (deseable),
        \emph{Won't} (no en esta entrega).
  \item \textbf{Modelo de Kano}: según la satisfacción del usuario (básicas, de rendimiento, de entusiasmo).
  \item \textbf{Criba de temas}: puntuar contra criterios (satisfacción, efectividad, ganar clientes, imagen,
        seguridad, productividad).
  \item \textbf{Peso relativo}: combina beneficio, penalización (por no tenerlo), coste y tamaño.
\end{itemize}

\subsection{Sprint Backlog (pila de Sprint) \diap{40}}
\begin{definicion}[Sprint Backlog]
Selección de requisitos del Product Backlog \textbf{negociados para el Sprint}: «el trabajo necesario del equipo para
alcanzar el objetivo del Sprint». Está \textbf{descompuesto en tareas} (lenguaje técnico) y es \textbf{propiedad del
equipo de desarrollo}.
\end{definicion}
\textbf{Dimensiones} \diap{58}: \textbf{prioridad} (primero las historias más prioritarias), \textbf{detalle} (cada
historia se divide en sus tareas) y \textbf{estado}: pendiente, en curso, terminada o \textbf{impedida}.

\subsection{Burndown Chart \diap{41--44}}
Representa el \textbf{trabajo pendiente} del equipo frente al tiempo. Hay dos tipos: del \textbf{Sprint} (días en el
eje X, horas o puntos pendientes en el Y) y del \textbf{proyecto} completo (Sprints en el eje X).
\begin{center}
\begin{tikzpicture}[font=\sffamily\footnotesize,x=0.95cm,y=0.05cm]
  \draw[->] (0,0) -- (10.6,0) node[right] {días};
  \draw[->] (0,0) -- (0,85) node[above] {pendiente};
  \foreach \x in {0,2,...,10} \node[below] at (\x,0) {\x};
  \draw[thick,gray,dashed] (0,80) -- (10,0) node[pos=0.08,above right,text=gray] {ideal};
  \draw[thick,mdaRed] (0,80) -- (1,79) -- (2,76) -- (3,74) -- (4,70) -- (5,64) node[above,text=mdaRed] {real (A)};
  \draw[thick,mdaGreen] (0,80) -- (1,68) -- (2,58) -- (3,47) -- (4,36) node[below left,text=mdaGreen] {real (B)};
\end{tikzpicture}
\end{center}
\begin{aviso}[cómo se interpreta \diap{44}]
\begin{itemize}[leftmargin=*]
  \item Línea real \textbf{por encima} de la ideal (A): el ritmo es \textbf{bajo}; conviene \textbf{reducir el alcance}
        del Sprint (renegociarlo con el PO).
  \item Línea real \textbf{por debajo} (B): el ritmo es \textbf{mayor} de lo previsto; quizá se acabe antes y se
        puedan \textbf{incorporar nuevas historias}.
\end{itemize}
Da información inmediata del Sprint y ayuda a preparar \textbf{acciones correctoras}. Si la línea se queda plana
varios días, hay un impedimento.
\end{aviso}

\subsection{Lista de impedimentos (Impediments Backlog) \diap{66--67}}
Sale del Daily: \textbf{repositorio de todo lo que impide alcanzar los objetivos del proyecto}, identificado y
priorizado. \textbf{Responsable: el Scrum Master.} Un impedimento puede ser la falta de una información o herramienta,
no poder reunirse con una persona clave, no tener un material, o \textbf{todo lo que amenace la calidad} del producto.
Cada impedimento se describe, se identifica su solución si es posible y se le \textbf{asigna un responsable}.

% ==============================================================================
\section{El Sprint y las cuatro reuniones}

\subsection{El ciclo de Scrum \diap{46}}
\begin{center}
\begin{tikzpicture}[font=\sffamily\footnotesize,node distance=0.5cm,
  b/.style={draw=mdaPrimary,fill=mdaLightAccent,rounded corners,align=center,minimum height=0.9cm,text width=2.3cm}]
  \node[b] (s0) {Sprint 0\\visión + backlog inicial};
  \node[b,right=0.6cm of s0] (pb) {Product\\Backlog};
  \node[b,right=0.6cm of pb] (pl) {Sprint\\Planning};
  \node[b,right=0.6cm of pl] (sb) {Sprint\\Backlog};
  \node[b,below=1.3cm of sb] (dd) {Sprint\\(1--4 semanas)\\+ Daily};
  \node[b,left=0.6cm of dd] (rv) {Sprint\\Review};
  \node[b,left=0.6cm of rv] (rt) {Sprint\\Retrospective};
  \node[b,fill=mdaLightGreen,draw=mdaGreen,left=0.6cm of rt] (inc) {Incremento\\terminado};
  \draw[-{Latex},thick] (s0)--(pb); \draw[-{Latex},thick] (pb)--(pl); \draw[-{Latex},thick] (pl)--(sb);
  \draw[-{Latex},thick] (sb)--(dd); \draw[-{Latex},thick] (dd)--(rv); \draw[-{Latex},thick] (rv)--(rt);
  \draw[-{Latex},thick,mdaGreen] (rv.south) |- ++(0,-0.5) -| (inc.south);
  \draw[-{Latex},thick,dashed,mdaPrimary] (rt.north) -- ++(0,0.55) -| (pl.south) node[pos=0.25,above] {siguiente Sprint};
\end{tikzpicture}
\end{center}

\subsection{El Sprint \diap{47--49}}
\begin{definicion}[Sprint]
Bloque de tiempo (\textbf{time-box}) de \textbf{1 a 4 semanas} durante el cual se crea un \textbf{incremento de producto
«terminado», utilizable y potencialmente desplegable}.
\end{definicion}
\begin{itemize}[leftmargin=*]
  \item La \textbf{longitud es fija} (mejor ritmo).
  \item Tiene una definición de qué se va a construir, un diseño y un plan \textbf{flexible}.
  \item \textbf{No se hacen cambios que afecten al objetivo del Sprint.}
  \item \textbf{Los objetivos de calidad no disminuyen.}
  \item El \textbf{alcance} sí se puede clarificar y renegociar entre el PO y el equipo a medida que se aprende.
\end{itemize}

\subsection{Duración de las reuniones (Sprint de 4 semanas) \diap{50}}
\begin{center}
\begin{tikzpicture}[font=\sffamily\footnotesize]
  \draw[rounded corners=12pt,thick,mdaPrimary] (0,0) rectangle (13,1.6);
  \node[draw,fill=mdaLightAccent,minimum width=2cm,minimum height=0.9cm,align=center] at (1.4,0.8) {Sprint\\Planning};
  \node[text=mdaRed] at (1.4,-0.3) {\textbf{8 horas}};
  \foreach \x in {3.3,4.3,5.3} \draw[mdaRed] (\x,0.8) circle (0.35) node {\scriptsize 1 día};
  \node at (6.6,0.8) {$\cdots$};
  \draw[mdaRed] (7.8,0.8) circle (0.35) node {\scriptsize 1 día};
  \node at (5.5,1.85) {Daily Scrum (cada día) \extra: 15 min};
  \node[draw,fill=mdaLightAccent,minimum width=1.9cm,minimum height=0.9cm,align=center] at (9.6,0.8) {Sprint\\Review};
  \node[text=mdaRed] at (9.6,-0.3) {\textbf{4 horas}};
  \node[draw,fill=mdaLightAccent,minimum width=2.1cm,minimum height=0.9cm,align=center] at (11.8,0.8) {Retrospective};
  \node[text=mdaRed] at (11.8,-0.3) {\textbf{3 horas}};
  \node[text=mdaRed] at (6.5,-0.9) {\textbf{Sprint de 4 semanas}};
\end{tikzpicture}
\end{center}
Para Sprints más cortos se reducen proporcionalmente. La duración del Daily (\textbf{15 minutos}) no aparece en la
diapositiva; es la de la Guía de Scrum.

\subsection{Sprint Planning \diap{51--59}}
\textbf{Reunión previa al comienzo de cada Sprint} para identificar el \textbf{objetivo del Sprint} y las tareas
concretas. Entrada: el Product Backlog revisado y priorizado. Salida: el \textbf{objetivo general del Sprint} y el
\textbf{Sprint Backlog}. Asegura que todo el equipo está alineado y comprometido con un plan claro y alcanzable.

\textbf{Objetivos} \diap{55}: definir el \textbf{objetivo del Sprint} (descripción corta y clara de lo que se quiere
lograr; guía el trabajo), \textbf{seleccionar los ítems} (por prioridad del PO y capacidad del equipo) y
\textbf{desarrollar un plan de trabajo} (descomponer en tareas y estimarlas).

\textbf{Dos fases} \diap{53}:
\begin{enumerate}[leftmargin=*]
  \item \textbf{Planificación (qué)}: se revisan los ítems del Product Backlog valorando su complejidad y se
        seleccionan, \textbf{siguiendo la prioridad}, los que caben en el Sprint.
  \item \textbf{Planificación detallada (cómo)}: el \textbf{equipo} traduce las historias seleccionadas al lenguaje de
        desarrollo y las divide en \textbf{tareas} (arquitectura, código, pruebas, marketing…).
\end{enumerate}
\begin{teorema}[selección de historias, paso a paso \diap{56}]
\begin{enumerate}[leftmargin=*,itemsep=1pt]
  \item El PO lee la primera HU del Product Backlog (la de \textbf{mayor prioridad}).
  \item La explica, detalla los \textbf{criterios de aceptación} y responde dudas.
  \item El equipo la \textbf{estima} (puntos de historia).
  \item Cada historia comprendida, valorada y seleccionada pasa al Sprint Backlog \textbf{manteniendo el orden de
        prioridad}.
  \item Cuando la suma de puntos alcanza (o supera ligeramente) la \textbf{velocidad media del equipo}, termina la
        primera fase.
\end{enumerate}
\end{teorema}
\textbf{Participantes} \diap{59}: el \textbf{PO} explica prioridades y negocia qué entra; el \textbf{equipo}
descompone en tareas y estima; el \textbf{SM} facilita y vela por el foco y la eficiencia.

\subsection{Backlog grooming (refinamiento) \diap{60--62}}
Revisar, actualizar y priorizar el Product Backlog para que esté \textbf{listo para los próximos Sprints}: depurar,
añadir historias, eliminar las que sobran, \textbf{dividir épicas y temas} y comprobar la priorización.
\begin{itemize}[leftmargin=*,itemsep=1pt]
  \item \textbf{No está formalmente incluida} entre las reuniones típicas de Scrum.
  \item Se hace \textbf{inmediatamente antes de cada planificación}.
  \item La \textbf{dirige el PO}; el equipo aporta detalle técnico, estimación y viabilidad; el SM facilita.
\end{itemize}

\subsection{Daily Meeting (Scrum diario) \diap{63--70}}
Reunión \textbf{diaria} con todo el equipo de desarrollo y el SM; \textbf{el PO puede asistir}. Cada miembro, de forma
voluntaria, responde:
\begin{center}
\begin{tcolorbox}[colback=mdaLightGreen,colframe=mdaGreen,width=0.8\textwidth,halign=center]
\textbf{1.} ¿Qué he hecho? \quad \textbf{2.} ¿Qué voy a hacer a continuación? \quad \textbf{3.} ¿Qué impedimentos tengo?
\end{tcolorbox}
\end{center}
\begin{itemize}[leftmargin=*]
  \item Breve y centrada en el objetivo. \textbf{No} se entra en detalle de lo hecho, \textbf{no} se discuten soluciones y
        \textbf{no} se divaga sobre el origen de un impedimento: solo cómo resolverlo y quién.
  \item Da a todos un conocimiento actualizado del estado; la información se refleja en un \textbf{panel} accesible.
  \item \textbf{Genera}: la lista de impedimentos (responsable el SM) y el Burndown actualizado.
  \item \textbf{Participantes} \diap{68}: el equipo (protagonista); el SM la facilita, \textbf{aunque participe no
        debería dirigirla}; el PO puede asistir, sin participar activamente, para aclaraciones rápidas.
  \item \textbf{Importancia} \diap{69--70}: transparencia y comunicación; adaptabilidad (ajustes diarios del plan);
        compromiso y responsabilidad (cada uno informa ante sus compañeros).
\end{itemize}
\begin{aviso}[el Daily no es]
Una reunión de estado para el jefe, una sesión para resolver problemas técnicos, ni una reunión donde el SM pregunta
y cada uno le contesta. Si surge un problema, se anota y se trata \emph{después} con quien haga falta.
\end{aviso}

\subsection{Sprint Review (revisión) \diap{71--79}}
\textbf{Objetivo}: recoger \emph{feedback} sobre el estado del producto. Es el punto de comunicación entre los
responsables del producto y el equipo. Tiene lugar \textbf{el último día del Sprint}. Se revisa \textbf{qué} se ha
hecho: puesta en común, novedades de la iteración y valor añadido.

\textbf{Finalidad} \diap{73}: recoger comentarios sobre el objetivo del Sprint, recoger la \textbf{aceptación} de lo
comprometido y analizar si hay que mejorar lo hecho o añadir ítems nuevos al Product Backlog.

\textbf{Preparación} \diap{74}: tener el material de la demo, un guion coherente y comprobar que todo funciona.

\textbf{Guion en 8 pasos} \diap{75--76}:
\begin{enumerate}[leftmargin=*,itemsep=1pt]
  \item Normas de la reunión si hay asistentes nuevos.
  \item Repaso de los objetivos del Sprint.
  \item El PO recapitula qué se ha terminado y qué no.
  \item El equipo revisa sus estadísticas y métricas, logros, problemas y cómo los resolvió.
  \item \textbf{Demostración} por el equipo de los ítems \textbf{terminados}, respondiendo dudas.
  \item El PO repasa el Product Backlog y el plan de entregas.
  \item Comentarios de stakeholders y PO sobre las novedades.
  \item Se recogen todos los comentarios para las siguientes reuniones de estimación y planificación.
\end{enumerate}
\textbf{Genera} \diap{77}: el \textbf{Product Backlog adaptado} a las nuevas necesidades y un \textbf{plan de entregas
actualizado} con la nueva velocidad y las nuevas prioridades.

\textbf{Participantes} \diap{78--79}: equipo (presenta y responde), SM (facilita), PO (voz del cliente: qué se ha
cumplido y prioridades futuras) y \textbf{stakeholders} (clientes, inversores, otros departamentos: su feedback guía el
producto). \textbf{Es la única reunión con stakeholders.}

\subsection{Sprint Retrospective \diap{80--91}}
\textbf{Después de la Review}, el equipo analiza \textbf{«el cómo»} trabaja para mejorar continuamente y ver qué le
impide ser más productivo.
\begin{itemize}[leftmargin=*]
  \item \textbf{Objetivo}: qué estamos haciendo bien, qué tenemos que mejorar (o dejar de hacer), y qué vamos a hacer en
        la siguiente iteración.
  \item \textbf{Características}: aprender de la experiencia, aumentar calidad y productividad, eliminar lo que no es
        útil y potenciar lo que sí, revisar riesgos y mejorar la comunicación.
  \item \textbf{Se revisa} \diap{84}: cada persona, las relaciones entre miembros, los procesos y las herramientas.
  \item \textbf{Participantes} \diap{85}: el equipo (todos participan), el SM (facilita, ambiente positivo y
        constructivo, foco en acciones de mejora) y el PO \textbf{no siempre es requerido}.
\end{itemize}
\textbf{Review frente a Retrospective}: la Review mira el \textbf{producto} (qué) con los stakeholders; la Retrospective
mira el \textbf{proceso} (cómo) con el equipo.

\subsubsection*{Técnicas de retrospectiva}
\begin{center}\small
\begin{tabular}{p{2.6cm}p{11.4cm}}
\toprule
\textbf{Técnica} & \textbf{Cómo funciona} \\ \midrule
\textbf{El Barco} \diap{86--87} & Se dibuja un barco que va hacia una isla. Post-its: \textbf{viento} (lo que nos impulsa), \textbf{ancla} (impedimentos que nos frenan), \textbf{iceberg} (riesgos), \textbf{sol} (lo que nos facilita la travesía) e \textbf{isla/tesoro} (el objetivo). \\
\textbf{Estrella de mar} \diap{88--89} & Cinco brazos: \textbf{Start doing} (probar algo nuevo), \textbf{Do more} (útil: hacerlo más y mejorarlo), \textbf{Continue to do} (útil y ya refinado: mantener), \textbf{Do less} (aporta poco), \textbf{Stop doing} (no aporta valor: eliminar). \\
\textbf{Las 4 L} \diap{90--91} & Cuatro cuadrantes: \textbf{Liked} (qué gustó), \textbf{Lacked} (qué faltó), \textbf{Learned} (qué aprendimos), \textbf{Longed for} (qué anhelamos). \\
\bottomrule
\end{tabular}
\end{center}

\subsection{Tabla resumen de las reuniones}
\begin{center}\footnotesize
\setlength{\tabcolsep}{3pt}
\begin{tabular}{p{2.1cm}p{2.3cm}p{1.6cm}p{3.1cm}p{2.7cm}p{2.6cm}}
\toprule
\textbf{Reunión} & \textbf{Cuándo} & \textbf{Duración (4 sem.)} & \textbf{Quién} & \textbf{Para qué} & \textbf{Genera} \\ \midrule
Sprint Planning & antes de cada Sprint & 8 h & PO, equipo, SM (facilita) & qué y cómo se hará & objetivo del Sprint y Sprint Backlog \\
Daily & cada día & 15 min \extra & equipo y SM; PO puede asistir & sincronizarse; detectar impedimentos & impedimentos, Burndown \\
Sprint Review & último día del Sprint & 4 h & equipo, SM, PO, \textbf{stakeholders} & inspeccionar el \textbf{producto} (qué) & PB adaptado, plan de entregas \\
Retrospective & después de la Review & 3 h & equipo y SM; PO opcional & mejorar el \textbf{proceso} (cómo) & acciones de mejora \\
(Grooming) & justo antes de cada Planning & -- & PO (dirige), equipo, SM & dejar el PB listo & PB refinado \\
\bottomrule
\end{tabular}
\end{center}

% ==============================================================================
\section{Sprint 0 y visión del producto (Elevator Pitch)}

\subsection{Sprint 0 \diap{93--95}}
Etapa en la que se define \textbf{la misión del trabajo, las herramientas y el equipo}. Es muy importante y de
\textbf{duración variable pero no indefinida}.
\begin{itemize}[leftmargin=*]
  \item \textbf{Viabilidad}: al final del Sprint 0 debe quedar claro si el proyecto es \textbf{viable} y tiene los medios
        y el apoyo necesarios.
  \item \textbf{Contrato}: el PO consigue apoyos y recursos, \textbf{selecciona al equipo y al Scrum Master}, acuerda
        \textbf{alcance y fechas} con el cliente y fija las condiciones para llegar a buen fin.
\end{itemize}
\begin{center}\small
\begin{tabular}{lp{11cm}}
\toprule
\textbf{Pregunta} & \textbf{Actividades del Sprint 0} \\ \midrule
¿Qué hacer? & Desarrollar la \textbf{visión del producto} (alcance) y el \textbf{Product Backlog inicial}. \\
¿Cómo hacerlo? & \textbf{Organizar el equipo} y la forma de trabajar (mecánica, herramientas de comunicación, \textbf{duración de los Sprints}, estándares, reglas); \textbf{análisis inicial} de la solución (arquitectura software inicial); \textbf{prueba de concepto}. \\
¿Cuándo hacerlo? & Desarrollar un \textbf{plan de entregas}. \\
\bottomrule
\end{tabular}
\end{center}

\subsection{Visión del producto \diap{96--101}}
\begin{definicion}[Visión]
Responde a \textbf{dónde se quiere estar} o \textbf{qué se quiere tener}. La define el \textbf{Product Owner}, y debe ser
\textbf{conocida, comprendida y compartida por todo el equipo}. Es el \textbf{origen y el final} del proceso: lo primero
que se hace y el punto al que se quiere llegar.
\end{definicion}
\textbf{Características} \diap{98}: información de \textbf{alto nivel} y propósito general; \textbf{sencilla, clara,
concisa y sin ambigüedades}; contiene la \textbf{esencia} de lo que se crea y sus claves; \textbf{fácilmente
comunicable}; instrumento vital para \textbf{conseguir presupuesto}; material de apoyo para el \textbf{lanzamiento,
venta y promoción}.

\textbf{Elementos básicos} \diap{99--100}:
\begin{multicols}{2}\small
\textbf{Concepto} (cómo se benefician los \emph{usuarios}):
\begin{itemize}[leftmargin=*,itemsep=1pt]
  \item ¿Qué problema se intenta resolver?
  \item ¿Cómo es la solución general?
  \item ¿Quiénes son los usuarios o clientes?
  \item ¿En qué situación están?
  \item ¿Qué necesidad suya se cubre?
  \item ¿Quiénes son los competidores y cómo nos comparamos?
\end{itemize}
\columnbreak
\textbf{Modelo de negocio} (cómo se benefician los \emph{creadores}):
\begin{itemize}[leftmargin=*,itemsep=1pt]
  \item ¿Cómo se consigue el presupuesto?
  \item ¿Cuál es el presupuesto y el tiempo para crear y lanzar el producto?
  \item ¿Cómo se genera beneficio?
  \item ¿Cómo encuentran y consiguen el producto los usuarios?
  \item ¿Cuáles son los costes del producto?
\end{itemize}
\end{multicols}

\subsection{Test del ascensor \diap{101}}
\begin{teorema}[Geoffrey Moore]
«La visión debería explicarse durante el tiempo que se tarda en subir o bajar varios pisos en un ascensor»: unos
\textbf{30 segundos a 2 minutos}. Si no cabe, no es lo bastante clara.
\end{teorema}
\textbf{Plantilla de Moore} (\emph{Crossing the Chasm}) \extra, la forma habitual de escribir el \emph{elevator pitch}:
\begin{tcolorbox}[colback=mdaYellowCard,colframe=mdaYellowBorder]
\textbf{Para} [cliente objetivo] \textbf{que} [necesidad u oportunidad], \textbf{el} [nombre del producto] \textbf{es
un} [categoría de producto] \textbf{que} [beneficio clave]. \textbf{A diferencia de} [competencia principal],
\textbf{nuestro producto} [diferencia principal].
\end{tcolorbox}

% ==============================================================================
\section{Trampas típicas}
\begin{center}\small
\begin{tabular}{p{6.6cm}p{7.4cm}}
\toprule
\textbf{Falso} & \textbf{Lo correcto} \\ \midrule
El Scrum Master asigna tareas al equipo. & El equipo se auto-organiza; el SM facilita y elimina impedimentos. \\
El equipo prioriza el Product Backlog. & Lo prioriza el \textbf{PO}; el equipo \textbf{estima}. \\
El PO es dueño del Sprint Backlog. & El Sprint Backlog es del \textbf{equipo}; el Product Backlog, del PO. \\
El Product Backlog está cerrado tras el Sprint 0. & Es \textbf{emergente} y dinámico: cambia siempre. \\
Durante el Sprint se puede cambiar su objetivo. & No; solo se \textbf{clarifica o renegocia el alcance}. \\
Si vamos mal, bajamos la calidad. & Los objetivos de calidad \textbf{no disminuyen}; se reduce el alcance. \\
En el Daily se resuelven los problemas técnicos. & Solo se identifican impedimentos; se resuelven después. \\
El SM dirige el Daily. & Lo facilita, pero no debería dirigirlo. \\
La Retrospectiva revisa el producto. & Revisa el \textbf{proceso} (cómo); la Review, el producto (qué). \\
Los stakeholders van a la Retrospectiva. & Van a la \textbf{Review}; la Retro es del equipo y el SM. \\
El grooming es una de las cuatro reuniones. & No está formalmente incluida; se hace justo antes de cada Planning. \\
Los Sprints varían de duración según el trabajo. & Duración \textbf{fija}, de 1 a 4 semanas. \\
Un rol es un puesto en la empresa. & Es una responsabilidad en el proceso. \\
El equipo tiene un subequipo de pruebas. & Scrum no reconoce sub-equipos; es multifuncional. \\
El Sprint 0 puede durar lo que haga falta. & Duración variable pero \textbf{no indefinida}. \\
\bottomrule
\end{tabular}
\end{center}

% ==============================================================================
\section{El examen en grupo: cómo organizaros}
Es una hora, sin apuntes y en grupo. La ventaja es que cada uno puede llevar dominado un bloque; el riesgo es perder
tiempo discutiendo.
\begin{enumerate}[leftmargin=*]
  \item \textbf{Repartid el temario antes del viernes}, con solapes: cada bloque, al menos dos personas.
        \begin{itemize}[itemsep=1pt]
          \item Historia, definición, pilares y valores, roles.
          \item Artefactos (Product Backlog DEEP, priorización, Sprint Backlog, Burndown, impedimentos).
          \item Reuniones (Planning con sus fases y pasos, grooming, Daily, Review con su guion, Retrospectiva y técnicas).
          \item Sprint 0, visión y elevator pitch.
        \end{itemize}
  \item \textbf{En el examen}: 5 minutos para leer todo y repartir preguntas por bloques; los ejercicios largos
        (visión, backlog, planificación) los hacen dos personas juntas; dejad 10 minutos para que otro revise cada
        respuesta.
  \item \textbf{Un único redactor por pregunta}, con letra clara. Los demás dictan y comprueban.
  \item Las \textbf{preguntas cortas} se contestan con la definición y un ejemplo; los \textbf{ejercicios}, con el
        resultado y el porqué (en qué fase, qué rol, qué artefacto).
\end{enumerate}

% ==============================================================================
\section{Ejercicios resueltos}

\begin{ejercicio}[1. Visión y elevator pitch]
Una start-up quiere crear \emph{ComparteCoche UGR}, una app para que estudiantes de la UGR que viven en pueblos del
área metropolitana compartan coche hasta los campus. Redacta la visión del producto con sus dos elementos básicos y
un \emph{elevator pitch}.
\end{ejercicio}
\textbf{Concepto.} Problema: muchos estudiantes de pueblos tienen malas conexiones de autobús y van solos en coche,
con gasto alto y problemas de aparcamiento. Solución: app que empareja conductores y pasajeros por horario de clases y
campus. Usuarios: estudiantes con y sin coche. Necesidad: llegar a tiempo, barato. Competencia: BlaBlaCar (pensado
para viajes largos, no diarios) y grupos de WhatsApp (sin verificación ni reparto de gastos).\\
\textbf{Modelo de negocio.} Presupuesto: subvención de la UGR y del área de movilidad; 6 meses y 40.000~\euro{} para el
MVP. Beneficio: comisión de 0,20~\euro{} por trayecto y patrocinio de gasolineras. Distribución: tiendas de apps,
difusión en la web de la UGR y en las matrículas. Costes: desarrollo, servidores, verificación de identidad.\\
\textbf{Elevator pitch.} \emph{Para los estudiantes de la UGR que viven fuera de Granada y pierden tiempo y dinero cada
día en el trayecto, ComparteCoche UGR es una app de coche compartido que empareja automáticamente conductores y
pasajeros según su horario de clases. A diferencia de BlaBlaCar o de los grupos de WhatsApp, solo admite usuarios
verificados con correo de la UGR y reparte los gastos automáticamente.}

\begin{ejercicio}[2. Sprint 0]
Para el proyecto anterior, enumera qué haríais en el Sprint 0 organizado en «qué, cómo y cuándo», y qué tiene que estar
claro al acabar.
\end{ejercicio}
\textbf{Qué}: visión (ejercicio 1) y Product Backlog inicial con las primeras épicas e historias. \textbf{Cómo}:
equipo de 6 (PO de la start-up, SM, 5 desarrolladores multifuncionales); Sprints de 2 semanas; Jira y Slack; estándares
de código y Definition of Done; arquitectura inicial (app Flutter + API REST + base de datos); prueba de concepto del
emparejamiento por horario. \textbf{Cuándo}: plan de entregas (MVP en el Sprint 4, versión pública en el Sprint 8).
\textbf{Al final}: que el proyecto es \textbf{viable} y hay \textbf{contrato} (recursos, equipo, alcance y fechas
acordados con el cliente).

\begin{ejercicio}[3. Product Backlog DEEP y MoSCoW]
Escribe 6 historias para el Product Backlog, estímalas y priorízalas con MoSCoW. Justifica que es DEEP.
\end{ejercicio}
\begin{center}\small
\begin{tabular}{cp{9.6cm}cc}
\toprule
\textbf{Orden} & \textbf{Historia de usuario} & \textbf{PH} & \textbf{MoSCoW} \\ \midrule
1 & Como estudiante, quiero registrarme con mi correo de la UGR para que solo haya usuarios verificados. & 3 & Must \\
2 & Como conductor, quiero publicar un trayecto con hora y plazas para llevar pasajeros. & 5 & Must \\
3 & Como pasajero, quiero buscar trayectos por campus y hora para llegar a clase. & 5 & Must \\
4 & Como pasajero, quiero reservar plaza para asegurarme el viaje. & 3 & Should \\
5 & Como usuario, quiero valorar a los demás para viajar con gente de confianza. & 2 & Could \\
6 & Como conductor, quiero cobrar con tarjeta dentro de la app para no manejar efectivo. & 8 & Won't (esta entrega) \\
\bottomrule
\end{tabular}
\end{center}
\textbf{DEEP}: \emph{detallado} (las de arriba tienen criterios de aceptación; la 6 puede quedarse como épica),
\emph{emergente} (en la Review pueden aparecer más), \emph{estimado} (puntos de historia por el equipo) y
\emph{priorizado} (orden por valor, lo decide el PO).

\begin{ejercicio}[4. Sprint Planning con velocidad]
La velocidad media del equipo es de 12 puntos. El Product Backlog priorizado es: A (5), B (3), C (8), D (2), E (3).
¿Qué historias entran en el Sprint según el procedimiento de las diapositivas? ¿Y si el PO dice que D es imprescindible
para el objetivo del Sprint?
\end{ejercicio}
Se recorren en \textbf{orden de prioridad} hasta alcanzar (o superar ligeramente) la velocidad: A (5) $\to$ 5; B (3) $\to$
8; C (8) $\to$ 16, que supera bastante los 12. Lo normal es parar en A y B (8 puntos) y negociar con el PO: si no se
puede partir C, se puede añadir D (2) $\to$ 10, y E no cabe entera (13 superaría ligeramente: aceptable si el equipo se
compromete). El orden de prioridad se mantiene en el Sprint Backlog. Si D es imprescindible para el objetivo, el PO
debe \textbf{subirla de prioridad} en el Product Backlog antes de seleccionar; no la elige el equipo por su cuenta.

\begin{ejercicio}[5. De historia a tareas (Sprint Backlog)]
Descompón la historia «Como pasajero, quiero buscar trayectos por campus y hora» en tareas y di qué estados pueden
tener.
\end{ejercicio}
Tareas (lenguaje técnico): diseñar la pantalla de búsqueda (3 h); servicio de búsqueda en la API, \texttt{GET /trayectos} filtrando por campus y hora (5 h);
índice en la base de datos por campus y hora (2 h); filtro de plazas libres (2 h); pruebas de aceptación de los
criterios (3 h); pruebas en dispositivo (2 h). Estados: \textbf{pendiente, en curso, terminada e impedida}. Las hace el
\textbf{equipo} en la segunda fase del Sprint Planning; el Sprint Backlog es suyo.

\begin{ejercicio}[6. Interpretar un Burndown]
Sprint de 10 días con 80 horas. El día 5 quedan 64 horas. ¿Qué indica y qué haríais? ¿Y si quedaran 30?
\end{ejercicio}
La línea ideal el día 5 estaría en 40 horas. Con 64 la real va \textbf{por encima}: el ritmo es bajo. Hay que buscar
la causa en el Daily (¿impedimentos?) y, si no se recupera, \textbf{reducir el alcance} renegociándolo con el PO, sin
bajar la calidad ni cambiar el objetivo del Sprint. Con 30 va \textbf{por debajo}: el ritmo es mayor y quizá se puedan
\textbf{incorporar nuevas historias} (las siguientes en prioridad, de acuerdo con el PO).

\begin{ejercicio}[7. ¿Qué está mal? (roles y reuniones)]
Señala el error en cada situación: (a) El Scrum Master reparte las tareas del Sprint en el Daily. (b) A mitad de Sprint,
el director comercial pide al equipo una funcionalidad nueva urgente y el equipo la mete. (c) La Retrospectiva la
dirige el cliente para ver la demo. (d) Como vamos retrasados, nos saltamos las pruebas.
\end{ejercicio}
(a) El equipo se auto-organiza; el SM facilita el Daily sin dirigirlo ni asignar trabajo. (b) Las peticiones van al
\textbf{PO}, que las añade y prioriza en el Product Backlog; el SM \textbf{protege} al equipo de interferencias y no se
cambia el objetivo del Sprint. (c) La demo con el cliente es la \textbf{Review}; la Retrospectiva es del equipo para
mejorar su proceso. (d) Los objetivos de calidad \textbf{no disminuyen}: se reduce el alcance.

\begin{ejercicio}[8. Preparar una retrospectiva]
El equipo ha tenido un Sprint con muchos problemas de despliegue y comunicación con el PO, pero ha aprendido a usar
pruebas automáticas. Plantea una retrospectiva con la estrella de mar.
\end{ejercicio}
\textbf{Start doing}: despliegue automático con integración continua. \textbf{Do more}: pruebas automáticas (útiles,
seguir ampliándolas). \textbf{Continue to do}: el panel de tareas actualizado tras el Daily. \textbf{Do less}:
reuniones largas por chat que no aportan. \textbf{Stop doing}: desplegar a mano el último día. Salida: 2 o 3 acciones
concretas para el siguiente Sprint, con responsable (por ejemplo, «sesión semanal de 30 minutos con el PO»).

% ==============================================================================
\section{Preguntas cortas con respuesta modelo}
\begin{enumerate}[leftmargin=*,itemsep=4pt]
  \item \textbf{¿Qué es Scrum?} Un marco de trabajo iterativo e incremental que no sigue un plan fijo, sino que se adapta
        continuamente a la evolución del proyecto.
  \item \textbf{¿Por qué «rugby»?} Takeuchi y Nonaka oponen el relevo (fases secuenciales) al rugby (el equipo avanza
        unido pasándose el balón): los equipos Scrum hacen un poco de todo, todo el tiempo.
  \item \textbf{Pilares empíricos}: transparencia, inspección y adaptación.
  \item \textbf{Elementos de Scrum}: 3 artefactos (Product Backlog, Sprint Backlog, Burndown), 4 reuniones
        (Planning, Daily, Review, Retrospective), 3 roles (PO, SM, equipo) y las reglas de trabajo.
  \item \textbf{Roles gallina y cerdo}: cerdos, el equipo Scrum (PO, SM, equipo), comprometidos; gallinas, los
        stakeholders (usuario, cliente, dirección), interesados pero no comprometidos.
  \item \textbf{Responsabilidades del PO}: visión compartida, requisitos, Product Backlog priorizado, aceptar o rechazar
        el incremento, resolver dudas, plan de entregas y ROI.
  \item \textbf{¿Por qué el SM no es un jefe?} Es facilitador, coach y mentor: elimina impedimentos y protege al equipo,
        pero no asigna tareas ni decide el producto.
  \item \textbf{Características del equipo}: auto-organizado, multidisciplinar, multifuncional, comprometido, sin
        sub-equipos, a tiempo completo, de 5 a 9 personas.
  \item \textbf{Lista DEEP}: detallada, emergente, estimada y priorizada.
  \item \textbf{Product Backlog frente a Sprint Backlog}: el primero es todo lo que podría necesitar el producto, en
        lenguaje de negocio, del PO; el segundo, lo seleccionado para el Sprint, descompuesto en tareas, del equipo.
  \item \textbf{Fases del Sprint Planning}: planificación (qué historias entran, por prioridad y velocidad) y
        planificación detallada (cómo: división en tareas por el equipo).
  \item \textbf{Las tres preguntas del Daily}: qué he hecho, qué voy a hacer y qué impedimentos tengo.
  \item \textbf{¿Qué genera la Review?} Product Backlog adaptado y plan de entregas actualizado con la nueva velocidad.
  \item \textbf{Review frente a Retrospective}: producto (qué, con stakeholders) frente a proceso (cómo, el equipo).
  \item \textbf{¿Qué es el Sprint 0?} Etapa de duración variable pero no indefinida para definir la visión, el backlog
        inicial, el equipo y la forma de trabajar; al final se sabe si el proyecto es viable.
  \item \textbf{Elementos de la visión}: concepto (beneficio para los usuarios) y modelo de negocio (beneficio para los
        creadores).
  \item \textbf{Test del ascensor}: la visión tiene que poder explicarse en el tiempo de un viaje en ascensor (Moore).
  \item \textbf{¿Qué es el backlog grooming?} Refinar el Product Backlog justo antes de cada Planning; no es una reunión
        formal de Scrum; la dirige el PO.
\end{enumerate}

% ==============================================================================
\clearpage
\section*{Tabla resumen para el último repaso}
\addcontentsline{toc}{section}{Tabla resumen}
\begin{center}\small
\begin{tabular}{p{3.3cm}p{10.8cm}}
\toprule
Scrum & marco iterativo e incremental; adaptación continua, no seguir un plan \\
Historia & Takeuchi-Nonaka 1986 (rugby) · Sutherland, Easel 1993 · Schwaber, OOPSLA 1996 · libro Schwaber-Beedle 2001 · Scrum Alliance 2002 (Schwaber, Cohn) \\
Pilares & transparencia, inspección, adaptación \\
Elementos & 3 artefactos + 4 reuniones + 3 roles + reglas \\
PO & cliente/negocio; prioriza el PB; acepta o rechaza; ROI; plan de entregas \\
SM & facilitador, coach, mentor; impedimentos; escudo; proceso \\
Equipo & auto-organizado, multifuncional, sin sub-equipos, 5--9; estima; Sprint Backlog \\
Product Backlog & ordenado y priorizado, lenguaje de negocio, del PO; DEEP \\
Priorización & priority poker, MoSCoW, Kano, criba de temas, peso relativo \\
Sprint Backlog & historias del Sprint en tareas; del equipo; pendiente/en curso/terminada/impedida \\
Burndown & trabajo pendiente; por encima de la ideal: reducir alcance; por debajo: añadir historias \\
Sprint & time-box fijo de 1--4 semanas; incremento terminado y desplegable; no cambia el objetivo ni baja la calidad \\
Planning (8 h) & qué (selección por prioridad hasta la velocidad) + cómo (tareas) \\
Daily (15 min) & hecho / haré / impedimentos; SM facilita sin dirigir; PO puede asistir \\
Review (4 h) & último día; demo; stakeholders; PB adaptado + plan de entregas \\
Retro (3 h) & después de la Review; cómo trabajamos; Barco, Estrella de mar, 4 L \\
Grooming & no formal; justo antes del Planning; lo dirige el PO \\
Sprint 0 & visión + PB inicial; equipo y forma de trabajar; arquitectura; prueba de concepto; plan de entregas; viabilidad y contrato \\
Visión & del PO; concepto + modelo de negocio; test del ascensor (Moore) \\
\bottomrule
\end{tabular}
\end{center}

\end{document}
